\title{Lista de Exercícios: Números Aleatórios e Cifras de Bloco}
\author{Prof. Gabriel Rodrigues Caldas de Aquino}
\date{Compilado em: \\ \today}

\begin{document}

\maketitle

\section{Cifra de Bloco}

\begin{itemize}
    \item Explique, com suas palavras, os conceitos de \textbf{difusão} e \textbf{confusão} segundo Claude Shannon.
    \item Descreva a \textbf{estrutura de Feistel}. Demonstre por que o mesmo algoritmo de Feistel pode ser usado tanto para encriptação quanto para decriptação apenas invertendo a ordem das subchaves $K_i$. Mostre isso para uma rodada apenas.
\end{itemize}

\section{Calculando uma Rodada da Cifra de Feistel}

Considere uma cifra de Feistel com bloco de \textbf{8 bits}, dividida em duas metades de 4 bits cada: $L_0$ (bits mais significativos) e $R_0$ (bits menos significativos).
A função de rodada $F$ recebe 4 bits de entrada e uma subchave $K_i$, mas, por simplicidade, nesta questão definimos $F(x, K_i) = 1111$ (ou seja, sempre retorna os 4 bits ‘1’, independentemente de $x$ e da subchave).

\textbf{Utilize:}
\begin{itemize}
    \item Bloco: $10100110$
    \item $F(R_0, K_1) = 1111$
\end{itemize}

\textbf{Faça:}
\begin{enumerate}
    \item Calcule $L_1$ e $R_1$ (em binário).
    \item Concatene $L_1$ e $R_1$ para obter o bloco de 8 bits de saída após uma rodada e apresente-o também em binário.
    \item Realize o processo inverso (descriptografe o bloco) e mostre que o texto original é recuperado.
    \item Explique brevemente por que, com essa função $F$ constante, a cifra se torna insegura, indicando qual propriedade (confusão ou difusão) é prejudicada.
\end{enumerate}

\section{Hash simples}
Considere a função de hash iterativa simples definida por:
\begin{itemize}
    \item Recebe de entrada dados de quaisquer tamanhos (Para facilitar utilize dados com tamanho múltipos de 8bits)
    \item Inicialização $H_0=00000000$;
    \item Funcionamento: separa o dado em blocos de 8 bits
    \item Em cada bloco de 8 bits $B_i$, $H_{i}=H_{i-1}\oplus B_i$.
    \item Resutado: último bloco de 8 bits
\end{itemize}

\begin{enumerate}
    \item Mostre explicitamente (com um exemplo numérico) como é possível reordenar blocos sem alterar o hash.
    \item Proponha uma modificação simples (mantendo iteração por blocos) que torne esse esquema resistente à simples reordenação de blocos (não precisa torná-lo criptograficamente seguro, apenas eliminar a propriedade de comutatividade).


\end{enumerate}

\section{Propriedades de Hash}
\begin{enumerate}
    \item Que características são necessárias em uma função de hash segura?
    \item O que é resistência à resistência à pré-imagem?
    \item O que é resistência à segunda pré-imagem?
    \item O que é resistência à colisão forte?
    \item Quais são as implicações de não se possuir resistência a colisão forte?
    \item O que é a função de compactação de hash? O que torna uma função de compactação segura?
\end{enumerate}

\section{Colisões em Hash}
Uma função de hash produz um resumo de (m) bits.
\begin{enumerate}
    \item Usando o paradoxo do aniversário, estime o esforço (número de tentativas da ordem) necessário para encontrar uma colisão por ataque de força bruta.
    \item Para ($m=128$) e para ($m=256$), escreva a ordem de magnitude desse esforço e comente se cada tamanho é suficiente para segurança moderna frente a ataques com recursos massivos.
\end{enumerate}

\end{document}